% !TEX TS-program = xelatex
% !TEX encoding = UTF-8 Unicode
\documentclass[10pt,aspectratio=1610]{beamer}
%\documentclass[10pt]{beamer}

\usetheme[progressbar=frametitle,numbering=counter]{metropolis}
\usepackage{appendixnumberbeamer}


\usepackage{booktabs}
\usepackage[scale=2]{ccicons}

\usepackage{pgfplots}
\usepgfplotslibrary{dateplot}

\usepackage{xspace}
\usepackage{bm}

%\newtheorem{definition1}{Definition}

\newcounter{mydef}
\stepcounter{mydef}
\newenvironment{mydefinition}{\begin{block}{Definition \themydef\stepcounter{mydef}}}{\end{block}}

% Color definitions 
%\usepackage[dvipsnames]{xcolor}

% Remove "Figure #" from figures
\setbeamertemplate{caption}{\raggedright\insertcaption\par}

% Custom commands
% Metropolis theme name 
\newcommand{\themename}{\textbf{\textsc{metropolis}}\xspace}

% Derivatives
\newcommand{\pfrac}[3][]{\frac{\partial^{#1} #2}{\partial #3^{#1}}}
\newcommand{\dd}[3][]{\frac{\mathrm{d}^{#1} #2}{\mathrm{d} #3^{#1}}}
\newcommand{\pp}[2][]{\partial^{#1}_{#2}}


\newcounter{example}
\stepcounter{example}
\newcounter{exercise}
\stepcounter{exercise}

\newcommand{\exercisen}{\theexercise\stepcounter{exercise}}
\newcommand{\examplen}{\theexample\stepcounter{example}}

\renewcommand{\vec}{\mathbf}
\newcommand{\varone}{\vec{x}}
\newcommand{\argone}{(t,\varone)}
\newcommand{\R}{\mathbb{R}}
\newcommand{\CC}{\mathbb{C}}
\newcommand{\Q}{\mathbb{Q}}
\newcommand{\Z}{\mathbb{Z}}
\newcommand{\N}{\mathbb{N}}
\newcommand{\norm}[1]{\left\lVert#1\right\rVert}
\newcommand{\fone}{\vec{f}}
\newcommand{\ftwo}{\vec{g}}
\newcommand{\fthr}{\vec{h}}
\newcommand{\ffor}{\bm{\phi}}
\newcommand{\opone}{\mathcal{L}}
\newcommand{\optwo}{\mathcal{G}}
\renewcommand{\Re}{\operatorname{Re}}
\newcommand{\iprod}[2]{\langle #1,#2 \rangle}
\newcommand{\diff}{\mathrm{d}}
\newcommand{\fouriert}{\mathcal{F}}
\newcommand{\BigO}{\mathcal{O}}

\title{Finite differences}
\subtitle{Real space approximations}
\date{16/2/2021}
\date{}
%\author{18.303 Linear Partial Differential Equations: Analysis and Numerics}
\institute{18.303 Linear Partial Differential Equations: Analysis and Numerics}
\titlegraphic{\hfill\includegraphics[height=2em]{../MIT-logo.pdf}}

\begin{document}
	
	\maketitle
	
%	\begin{frame}{Table of contents}
%		\setbeamertemplate{section in toc}[sections numbered]
%		\tableofcontents%[hideallsubsections]
%	\end{frame}

%%%%%%%%%%%%%%%
%\section{Finite differences}
%%%%%%%%%%%%%%%

%--------------------
% Finite difference 

\begin{frame}{Finite differences}
	What is the simplest way to do differentials on a computer? 
		
	\pause
	For the sake of simplicity let's assume we have $ u:[0,1] \to \R $ and we want to make sense of $ \dd{u(x)}{x} $.
	
	\pause
	We \alert{discretize} the variable $ x $, the function $ u $, and the derivative $ u'(x) $. Formally this means that 
	\[ u(x) \to u_n = u(x_n), \quad u'(x) \to u'_n,\]
	where $ n=1,2,...,N. $
	
	\pause
	For now, let's assume that $ x_n = n \Delta x $, where $ \Delta x = 1/N $ (we will later talk about non-evenly spaced discretization). We can \emph{approximate} $ u'(x) $ by writing the Taylor series for $ u $:
	\[ u(x_n + \Delta x) = u(x_n) + u'(x_n) \Delta x + \BigO(\Delta x^2).  \]
 	(Here we use "Big-$ \BigO $ notation": the term $ \BigO (\Delta x^2) $ scales as $ \Delta x^2 $ approaching zero as $ \Delta x \to 0 $. Frankly it means that for small $ \Delta x $ this term is very small. ) 
\end{frame}

%--------------------
% Finite difference 2

\begin{frame}{Finite differences}
	We can solve this for $ u'(x_n) $ giving
	\begin{equation}\label{forward_taylor}
		u'(x_n) =  \frac{u(x_n+\Delta x) - u(x_n)}{\Delta x} + \BigO(\Delta x)
	\end{equation}
	i.e. 
	\[ u'_n \approx \frac{u_{n+1}-u_n}{\Delta x}. \]
	
	\pause
	We notice that this is just a linear combination of the terms in the vector $ u_n $ i.e. it can be expressed as a linear operation
	\[ \vec{u}' \approx \delta_+ \vec{u}. \]
	
	\pause
	We call the operator $ \delta_+ $ \alert{forward difference} operator. Changing the sign of $ \Delta x $ in Eq.~\eqref{forward_taylor} gives the \alert{backward difference} operation
	\[ u'_n \approx \frac{u_n - u_{n-1}}{\Delta x} = (\delta_- \vec{u})_n. \]
\end{frame}

%--------------------
% Finite difference 3

\begin{frame}{Finite differences}
	The fact that we have $ \BigO(\Delta x) $ in Eq.~\eqref{forward_taylor} tells that the difference operation is \emph{first order} i.e. the error of the operation when approximating the actual differential scales linearly with $ \Delta x $. 
	
	\pause
	We can do a bit better.  Calculating 
	\[ \vec{u}' \approx \left( \frac{\delta_+ + \delta_-}{2} \right)\vec{u} = \delta_0 \vec{u} \]
	gives in indices 
	\[ u'_n = \frac{u_{n+1}- u_{n-1}}{2 \Delta x} + \BigO(\Delta x^2). \]
	
	We call $ \delta_0 $ the \alert{central difference} operator and, as we can see, it's second order in $ \Delta x $. 
	
	
	%	Now, the function $ u $ is evaluated at points $ x_n $ and becomes a vector $ \vec{u} $. What about $ \delta $? We know that it's a linear operator from the vector space to itself. Since the vector space is now $ \R^N $, an operator that will do is a matrix $  $
\end{frame}


%--------------------
% Exercise

\begin{frame}{Exercise \exercisen}
	%	\metroset{block=fill}
%	\centering
%	\begin{block}{Inner product induces a norm}
		Derive a formula for the second order derivative $ u''_n $ by using the Taylor series for $ u_{n+1} $ and $ u_{n-1} $.
%	\end{block}	
\end{frame}

%--------------------
% Operators 1

\begin{frame}{Operators $ \delta $}
	Note that in the discrete setting our vector space is just $ \R^n $. What are the difference operators $ \delta $ then? % \pause
	They are linear maps i.e. matrices. 
	
	\pause
	To calculate a matrix for the second differential, let's write
	 \[ \delta^{(2)} \vec{u} = \frac{1}{\Delta x^2}\begin{pmatrix}
		u_0 - 2 u_1 + u_2 \\ u_1 - 2 u_2 + u_3 \\  \vdots \\ u_{N-1} - 2u_{N} + u_{N+1}
	\end{pmatrix}. \] 
	
	\pause
	We immediately notice that we need values $ u_0 $ and $ u_{N+1} $. In order to define the operator properly, we need a notion of boundary conditions.

\end{frame}

%--------------------
% Operators 2

\begin{frame}{Operators $ \delta $}
	We can write this e.g. for Dirichlet boundaries defining $ u_0 = u_{N+1} = 0 $. 
	
	\pause
	Now,
	\[ \delta^{(2)}_{\text{D}} \vec{u} = \frac{1}{\Delta x^2}\begin{pmatrix}
		 - 2 u_1 + u_2 \\ u_1 - 2 u_2 + u_3 \\  \vdots \\ u_{N-1} - 2u_{N}
	\end{pmatrix} \]
	and we can write
	\[ \delta^{(2)}_{\text{D}} = \frac{1}{\Delta x^2} \begin{pmatrix}
		-2 & 1  &   & &  &  \\
		1  & -2 & 1 &  &  & \\
		 & \ddots & \ddots & \ddots  & & \\
		 & & 1 & -2 & 1 & \\
		 & & & 1 & -2 &
	\end{pmatrix}. \]
	
	\pause
	This matrix takes vectors from $ \R^N $ to $ \R^N $ i.e. it's a square matrix. 
	
\end{frame}

%--------------------
% Poisson revisited

\begin{frame}{Poisson equation revisited}
	We can write our Poisson equation in the discretized setting as 
	\[ \delta^{(2)}_{\text{D}} \vec{u} = \vec{f}. \]
	
	This is a \emph{sparse} linear system that can be solved for any $ \vec{f} $ ($ \delta^{(2)}_{\text{D}} $ is invertible). 
	
	We can also formulate an eigenvalue problem 
	\[ \delta^{(2)}_{\text{D}} \ffor  = \lambda \ffor. \]
\end{frame}

%--------------------
% Exercise

\begin{frame}{Exercise \exercisen}
	Try the ansatz $ \phi_k^{(n)} = \sin(a_n k) $, $ k=1,...,N $ for solving the eigenvalue problem 
	\[ \delta^{(2)}_{\text{D}} \ffor^{(n)}  = \lambda_n \ffor^{(n)}. \]
	
%	What is the value of $ a_n $? What about $ \lambda_n $?
\end{frame}

%--------------------
% From polynomials

\begin{frame}{Finite differences from polynomial interpolation}
	Assume we are given three points $ 0 $, $ \Delta x $, and $ 2 \Delta x $ with corresponding values for $ u $: $ u_0 $, $ u_1 $, and $ u_2 $. 
	
	We can fit a second order polynomial 
	\[ p_2(x) = a_0 + a_1 x + a_2 x^2 \]
	by requiring that $ p_2(x_i) = u_i $. 
	
	\pause
	This gives a linear system 
	\begin{align*}
		u_{0}	&= p_2(0) = a_{0}, \\
		u_{1}	&= p_2(\Delta x) = a_{2}\Delta x^{2} + a_{1}\Delta x + a_{0}, \\
		u_{2}	&= p_2(2 \Delta x)=4a_{2}\Delta x^{2} + 2a_{1}\Delta x + a_{0}.
	\end{align*}
\end{frame}

\begin{frame}{Finite differences from polynomial interpolation}
	Solving for $ a_k $ gives 
	\begin{align*}
		a_{0}	&= u_{0}, \\
		a_{1}	&=\frac{-3 u_0 + 4 u_1 - u_2}{2 \Delta x}, \\
		a_{2}	&= \frac{u_0 - 2 u_1 + u_2}{2 \Delta x^2}.
	\end{align*}
	
	\pause
	We can evaluate 
	\[ p'(\Delta x) = \frac{u_2 - u_0}{2 \Delta x} \]
	and 
	\[ p''(\Delta x) = \frac{u_2 - 2 u_1 + u_0}{\Delta x^2}. \]
\end{frame}

\begin{frame}{Finite differences from polynomial interpolation}
	This gives us a systematic way of deriving stencils (operators) for finite differences on any number of points (we could have chosen e.g. 5 points and a 4th order polynomial). 
	
	\pause
	Moreover, we could have chosen our points with a non-equal spacing and could have still solved for the coefficients $ a_k $. This shows that this method can be used for non-uniform grids. 
	
	\pause
	The algorithm that automatically generates stencils from the interpolating polynomial forms is known as the Fornberg algorithm.
\end{frame}

\end{document}
